\section{Actions physiques et spectre de la mémoire réduite
}
\label{nuclear:12}
\subsection{Généalogie des sections dimensionnelles
}

Les coordonnées \(\pi,\varphi,e,\alpha\) sont les réalisations des sorties corrélées construites dans les sections~\ref{sec:alpha} et~\ref{sec:electron}. Dans cette section, on utilise \(\alpha=\alpha_A\), avec l’orientation et la périodisation de la définition~\ref{def:alpha-operativa} :




\begin{equation}
\alpha=\pi+e-\varphi-4-\kappa_{\rm per},\qquad
\kappa_{\rm per}=\frac{N_K}{1000^{12}-1},
\quad N_K=234543140729659824621058914794146601.
\label{nuc12:eq:2.1}
\end{equation}



L’identité~\eqref{eq:alpha-identidad-completa} se déduit de la récurrence avec retenue compatible et de son télescopage, démontrés dans la proposition~\ref{prop:alpha-telescopia}. La période douze correspond au registre : une fenêtre de douze blocs avec condition au bord nulle et une section infinie compatible sont des objets différents. L’égalité \eqref{nuc12:eq:2.1} utilise cette dernière. Le registre ordonné \(K\) conserve ses incidences tout en admettant la lecture scalaire \(\kappa_{\rm per}\).

Le lecteur d’action, construit dans la section~\ref{iv:action:section}, détermine




\begin{equation}
\begin{aligned}
H_5&=\tfrac12\sqrt{1000\alpha/\varphi}-\alpha+
\tfrac9{16}\alpha^2-\tfrac59\alpha^3+
\tfrac7{48}\alpha^4-\tfrac1{54}\alpha^5,\\
D_A&=\exp(-100\pi\alpha/9),\\
\mathcal C_\pi&=169\alpha^6+\frac{D_A\alpha^7}{1-D_A\alpha},\\
\eta_{\rm ret}&=H_5-\frac{90}{\pi}\mathcal C_\pi,\\
\hbar_{\rm pre}&=H_5\mathcal A_0,\qquad
\hbar_{\rm ret}=\eta_{\rm ret}\mathcal A_0,\qquad
R_{\rm act}=H_5/\eta_{\rm ret},
\end{aligned}
\label{nuc12:eq:2.2}
\end{equation}



Ici, \(\mathcal A_0=10^{-34}\mathcal U_S\). La décennie provient du lecteur déterminantal de \(\varphi\alpha^{16}\), dont la construction et l’encadrement figurent dans le théorème~\ref{iv:action:decada} ; \(H_5\) et \(\eta_{\rm ret}\) sont les mantisses des deux sections positives d’action réduite. Dans chaque section \(a\in\{\mathrm{pre},\mathrm{ret}\}\), on a \(h_a=2\pi\hbar_a\).

La construction de l’électron central, développée dans la section~\ref{sec:electron}, détermine la coordonnée




\begin{equation}
\begin{aligned}
\Delta_4&=\frac{\pi-e\log\pi}{270}\left(1+\frac{\pi}{729}\right),\\
B_e&=\frac{\sqrt3}{4}
\left[\exp(\varphi/\pi^2)-22\alpha^3\right](1+15\Delta_4),\\
\kappa_e&=80-54\alpha+6\Delta_4,\qquad
\varepsilon_e^\beta=B_eR_{\rm act}^{\kappa_e}.
\end{aligned}
\label{nuc12:eq:2.3}
\end{equation}



Le préfacteur \(\sqrt3/4\) procède de la norme du commutateur des projecteurs de somme et de produit sur les fibres spécifiées dans \eqref{eq:el-prefactor}. La route centrale et ses registres déterminent les coefficients de retour de \eqref{eq:el-kappa} et~\eqref{eq:el-beta}. La réalisation dimensionnelle est \(E_e^\beta=\varepsilon_e^\beta\mathcal U_E\), \(m_e^\beta=E_e^\beta/c^2\), avec le système de coordonnées énergétiques de \eqref{eq:el-dimensional}.

L’horloge dodécaphasique détermine la base énergétique de \eqref{eq:el-cuanto-reloj} :



\[
\omega_{108}=\frac{2\pi}{108t_0},\quad
E_{\rm clk}=\hbar_{\rm ret}\omega_{108},\quad
b_e=B_e\frac{\mathcal U_E}{E_{\rm clk}}.
\]



La substitution de \eqref{nuc12:eq:2.3}, en conservant la même énergie physique, démontre



\begin{equation}
\boxed{\frac{m_e^\beta c^2}{\hbar_{\rm ret}}
=b_eR_{\rm act}^{\kappa_e}\omega_{108},\qquad
\frac{m_e^\beta c^2}{\hbar_{\rm pre}}
=b_eR_{\rm act}^{\kappa_e-1}\omega_{108}.}
\label{nuc12:eq:2.4}
\end{equation}



La division par \(h_a\), au lieu de \(\hbar_a\), transforme la fréquence angulaire en fréquence cyclique. La conversion \(\mathcal U_E/E_{\rm clk}\) reste explicite : une coordonnée exprimée en \(\mathrm{MeV}\) change lorsqu’on remplace sa base énergétique par le quantum de l’horloge.

\subsection{Normalisation de l’action Maxwell--Dirac}

La réponse constitutive de la section~\ref{iii:sec:constitutiva} produit la vitesse \(c\), l’impédance \(Z\), la permittivité \(\varepsilon\) et la perméabilité \(\mu\), toutes positives, avec \(\varepsilon c=Z^{-1}\) et \(\mu c=Z\). La section positive de charge de \eqref{iii:eq:charge} satisfait



\begin{equation}
e_a^2=\frac{2\alpha h_a}{Z}
=4\pi\alpha\varepsilon\hbar_a c .
\label{nuc12:eq:3.1}
\end{equation}



Ici, \(e_a>0\) désigne la charge élémentaire de la section \(a\) ; \(e\) dans \eqref{nuc12:eq:2.1}--\eqref{nuc12:eq:2.3} désigne la coordonnée d’Euler. Nous fixons un système de coordonnées homogène avec \(c,Z,\varepsilon,\mu,\hbar_a,e_a\) des constantes positives, la masse \(m=m_e^\beta>0\) et la charge \(q\in\{-e_a,e_a\}\) ; alors \(\sigma=q/e_a\in\{-1,1\}\). Dans la réalisation de Minkowski de signature \((+,-,-,-)\), avec \(x^0=ct\), l’action Maxwell--Dirac est



\begin{equation}
\mathcal I=\int d^4x\left[
-\frac{F_{\mu\nu}F^{\mu\nu}}{4\mu c}
+\bar\psi\left(i\hbar_a\gamma^\mu
(\partial_\mu+iqA_\mu/\hbar_a)-mc\right)\psi
\right].
\label{nuc12:eq:3.2}
\end{equation}



La composante \(A_0\) est le potentiel scalaire divisé par \(c\), et \(F=dA\). La représentation spinorielle est fixée au moyen des matrices \(g^I\) de \eqref{vii:eq:clifford-material-explicito}, qui satisfont \(\{g^I,g^J\}=2\eta^{IJ}I_4\) pour \(\eta=\operatorname{diag}(-1,1,1,1)\). Nous définissons \(\Gamma^I=-ig^I\), de sorte que \(\{\Gamma^I,\Gamma^J\}=-2\eta^{IJ}I_4\), et, dans ce système de coordonnées plat, \(\gamma^\mu=\delta^\mu_I\Gamma^I\), \(\bar\psi=\psi^\dagger\Gamma^0\). Ainsi sont fixés conjointement la signature, l’adjoint et le facteur imaginaire ; les matrices \(\gamma^\mu\) sont distinguées du paramètre scalaire de Holst utilisé ensuite. On considère des champs lisses et des variations à support compact, ou des conditions aux limites qui annulent les termes de bord. L’action réalise les coefficients déjà engendrés et la représentation matérielle de la section~\ref{vii:subsec:accion-espinorial-afin}.

On définit



\[
\ell_a=\frac{\hbar_a}{mc},\quad X=x/\ell_a,\quad
\Psi=\ell_a^{3/2}\psi,\quad
\mathsf a_\mu=\frac{e_a\ell_a}{\hbar_a}A_\mu,\quad
\mathsf f=d\mathsf a.
\]



Alors



\begin{equation}
\boxed{\frac{\mathcal I}{\hbar_a}=
\int d^4X\left[
-\frac{\mathsf f_{\mu\nu}\mathsf f^{\mu\nu}}{16\pi\alpha}
+\bar\Psi\bigl(i\gamma^\mu(\partial_\mu+i\sigma\mathsf a_\mu)-1\bigr)\Psi
\right].}
\label{nuc12:eq:3.3}
\end{equation}



\begin{proof}
Le changement de volume apporte \(\ell_a^4\), chaque dérivée apporte \(\ell_a^{-1}\) et le produit des champs fermioniques apporte \(\ell_a^{-3}\). Le coefficient du terme de masse est \(mc\ell_a/\hbar_a=1\). Le terme électromagnétique apporte \(\hbar_a/(4\mu c e_a^2)\). Comme \(\mu c=Z\) et \(Ze_a^2=4\pi\alpha\hbar_a\), ce coefficient est \(1/(16\pi\alpha)\). Les deux termes et leurs normalisations sont déterminés.
\end{proof}

La variation par rapport à \(\mathsf a_\nu\), suivie de la substitution de \eqref{nuc12:eq:2.1}, donne




\begin{equation}
\partial_\mu\mathsf f^{\mu\nu}
=4\pi\left(\pi+e-\varphi-4-\kappa_{\rm per}\right)
\sigma\bar\Psi\gamma^\nu\Psi .
\label{nuc12:eq:3.4}
\end{equation}



L’annulation compose les sections d’action, la construction électronique et la réponse constitutive ; le lecteur régional et dodécaphasique détermine le couplage restant. La masse persiste dans la longueur \(\ell_a\), et l’action dans la conversion dimensionnelle. La normalisation de Compton est une réalisation ultérieure de ces sorties. Dans la réalisation coulombienne du même électron, la longueur de Bohr est \(\ell_a/\alpha\) : substituer \(e_a^2=4\pi\alpha\varepsilon\hbar_a c\) dans \(4\pi\varepsilon\hbar_a^2/(m e_a^2)\) donne exactement cette expression.


La déduction utilise des coefficients homogènes fixes. Pour des coefficients dépendant de l’espace-temps, la règle du produit introduit les dérivées des échelles et du couplage dans les équations transformées. La variation d’une action à paramètres variables doit conserver ces termes.

\subsection{Échelle surfacique de l’action gravitationnelle et de l’entropie}

La réalisation circulaire du cycle de \(108\) pas, construite dans la section~\ref{v:ant:sec:gravity}, transporte la longueur \(L=L_*\alpha^{16}5759/23040\) de~\eqref{xii:radial:length}. Avec \(t_0=\pi L/(54c)\), elle détermine



\begin{equation}
t_P=\frac{54}{\pi}t_0,\qquad
G_a=\frac{c^5t_P^2}{\hbar_a},\qquad
\ell_P^2=\frac{\hbar_aG_a}{c^3}=c^2t_P^2.
\label{nuc12:eq:4.1}
\end{equation}



La condition de coïncidence des rayons de cette réalisation est \(G_a m_{108,a}/c^2=R_{108}=L\), avec \(m_{108,a}=\hbar_a\omega_{108}/c^2\). C’est la restriction au mode unité de~\eqref{xii:radial:G} ; la forme chronologique conserve le facteur \((54/\pi)^2\). Le lecteur emploie le rayon réduit \(G_a m/c^2\). La réciprocité \(G_{\rm pre}/G_{\rm ret}=R_{\rm act}^{-1}\) conserve \(\ell_P^2=L^2\) entre les deux sections. La référence \(L_*\), la base d’arête \(u_L\) et l’avancée \(ct_0\) sont reliées par~\eqref{xii:radial:clock}.

La section~\ref{vii:sec:realizacion-cartan} construit la cotétrade et la connexion à partir du transport déjà réalisé. Dans l’action de la section~\ref{vii:sec:corriente-respuesta-cartan}, nous fixons une cotétrade non dégénérée \(e^I\), l’orientation, la signature \(\eta=\operatorname{diag}(-1,1,1,1)\) et une connexion métrique \(\omega\), de courbure \(F^{IJ}=d\omega^{IJ}+\omega^I{}_{K}\wedge\omega^{KJ}\). Nous prenons \(\gamma\in\mathbb R\setminus\{0\}\) et maintenons \(\gamma,\Lambda,c,G_a,\hbar_a\) constants dans la variation. Nous posons \(\Sigma^{IJ}=e^I\wedge e^J\) et \(\mathsf P_\gamma=\star+\gamma^{-1}I\) ; la dualité interne lorentzienne satisfait \(\star^2=-I\) sur les bivecteurs. Les variations ont un support compact ou des conditions aux limites qui annulent le terme de surface. La substitution de \eqref{nuc12:eq:4.1} dans \eqref{vii:eq:accion-cartan-holst} donne




\begin{equation}
\boxed{\frac{\mathcal I_{\rm grav}}{\hbar_a}
=\frac1{16\pi\ell_P^2}
\left[\int\Sigma_{IJ}\wedge(\mathsf P_\gamma F)^{IJ}
-2\Lambda\int\operatorname{vol}_e\right].}
\label{nuc12:eq:4.2}
\end{equation}



La cotétrade et les champs matériels étant fixes pendant la variation de la connexion, soit \(\sigma_{IJ}\) le courant matériel défini par
\(\delta_\omega\mathcal I_m=-\tfrac12\int\delta\omega^{IJ}\wedge\sigma_{IJ}\), selon \eqref{vii:eq:corriente-variacional}. La même substitution dans le théorème~\ref{vii:thm:ecuacion-cartan-holst} produit



\begin{equation}
D_\omega(\mathsf P_\gamma\Sigma)
=8\pi\ell_P^2\,\frac{\sigma}{\hbar_a}.
\label{nuc12:eq:4.3}
\end{equation}



\begin{proof}
Dans la section \(a\), le coefficient original est \(1/(2\kappa_a c)\), avec \(\kappa_a=8\pi G_a/c^4\). La division par \(\hbar_a\) le transforme en \(1/(16\pi\ell_P^2)\) ; le terme cosmologique conserve son facteur deux. Dans l’équation de connexion, on a \(\kappa_a c\hbar_a=8\pi\ell_P^2\). Les deux identités utilisent la même section dimensionnelle.
\end{proof}

La fonctionnelle de Barbero de la section~\ref{v:ant:sec:barbero} conserve son expression intégrale. Ses canaux sont \(q_\pm=e^{-x\mp y}\), avec \(x=\pi A/180\), \(y=\pi C^*/180\), où \(A\) et \(C^*\) sont les coordonnées angulaires exprimées en degrés. Le domaine \(x>|y|\) assure \(0<q_\pm<1\) :



\begin{equation}
\gamma=\frac{180}{\pi}xy+
12[S_{90}(q_-)-S_{90}(q_+)]
-[S_{120}(q_-)-S_{120}(q_+)],\qquad
S_m(q)=\frac{q^m}{1-q^{3m}}.
\label{nuc12:eq:4.4}
\end{equation}



L’incidence hexadique et octadique fixe les multiplicités \(90\) et \(120\), et la chaîne orientée fixe les coefficients \(12\) et \(-1\), comme le démontre \eqref{v:ant:eq:gamma}. La construction aréale de la section~\ref{v:ant:sec:area} donne en outre



\begin{equation}
a_0=\frac{1+2\cos(\pi/18)}3\,\frac y6,\qquad
\widehat{\mathcal A}=\gamma a_0\ell_P^2(N_{90}+N_{120}).
\label{nuc12:eq:4.5}
\end{equation}



L’inversion \(y\mapsto-y\) change les signes de \(\gamma\) et de \(a_0\) ; leur produit est pair. L’opérateur \(\mathsf P_\gamma\) conserve, en revanche, sa dépendance à l’orientation. L’invariance du produit aréal se distingue ainsi de l’invariance de chaque coefficient de l’opérateur. Le spectre aréal se réalise dans la section positive \(\gamma a_0>0\), avec les opérateurs d’occupation \(N_{90},N_{120}\geq0\) construits dans la section~\ref{v:ant:sec:area}.


Pour l’horizon cosmologique sphérique de \eqref{vii:eq:horizonte-cosmologico}, de rayon \(R>0\) et d’aire \(\mathcal A_H=4\pi R^2\), nous fixons une section \(a\) et nous abrégeons \(\hbar=\hbar_a\), \(G=G_a\) pendant ce calcul et celui du rebond qui suit. Alors



\begin{equation}
\frac{S_H}{k_B}=\frac{\mathcal A_H}{4\ell_P^2},\qquad
T_{\rm cos}=\frac{\hbar c}{2\pi k_BR},\qquad
E_\Lambda=\frac{c^4R}{2G},\qquad T_{\rm cos}S_H=E_\Lambda .
\label{nuc12:eq:4.6}
\end{equation}



L’échelle surfacique de \eqref{nuc12:eq:4.2}, \eqref{nuc12:eq:4.3} et \eqref{nuc12:eq:4.6} est la même. La température correspond à l’horizon cosmologique sphérique ; la normalisation de l’horizon de Schwarzschild constitue une autre réalisation. En faisant varier \(R\) tout en maintenant \(c,G,\hbar,k_B\) fixes, la différentiation de \eqref{nuc12:eq:4.6} conserve les deux termes : \(T_{\rm cos}\,dS_H=2\,dE_\Lambda\) et \(S_H\,dT_{\rm cos}=-\,dE_\Lambda\), comme dans \eqref{vii:eq:diferencial-horizonte}. On obtient la différentielle d’une identité entre les états de cette collection, et non un remplacement du bilan dynamique avec d’autres paramètres variables.

\subsubsection{Une application dynamique explicite}

Le théorème~\ref{vii:thm:amplitud-material-explicita} détermine, pour un état spinoriel à second moment conservé \(q_\varrho>0\), l’amplitude suivante. La conservation s’exprime par \eqref{vii:eq:conservacion-momento-cuartico} ; le corollaire~\ref{vii:cor:dominio-positivo-conservado} fournit une réalisation concrète dans le secteur à trois occupations :



\begin{equation}
s_0=\frac{3\pi G\hbar^2}{2c^2V_c^2}
\frac{\gamma^2}{1+\gamma^2}q_\varrho .
\label{nuc12:eq:4.7}
\end{equation}



Dans la réalisation spatialement plate avec une matière de pression nulle et \(\Lambda=0\), la constante \(\rho_0>0\) est une densité d’énergie, \(V_c>0\) est le volume comobile de référence et \(E_0=\rho_0V_c>0\). La solution du théorème~\ref{vii:thm:rebote-polvo} est




\[
a_b^3=s_0/\rho_0,\qquad
\tau^2=\frac{s_0c^2}{6\pi G\rho_0^2},\qquad
a(t)=a_b\left(1+\frac{(t-t_b)^2}{\tau^2}\right)^{1/3}.
\]



La restitution de \(c\) dans \eqref{vii:eq:dinamica-homogenea} utilise le coefficient \(8\pi G/(3c^2)\) lorsque la densité est énergétique. En particulier, \((\dot a/a)^2=(8\pi G/(3c^2))(\rho_0a^{-3}-s_0a^{-6})\). Substituer \eqref{nuc12:eq:4.7} dans les paramètres de la solution démontre



\begin{equation}
\boxed{\left(\frac{\tau E_0}{\hbar}\right)^2
=\frac{q_\varrho}{4}\frac{\gamma^2}{1+\gamma^2},\qquad
\mathcal V_b=\frac{3\pi}{2}q_\varrho
\frac{\gamma^2}{1+\gamma^2}\ell_P^2\frac{\hbar c}{E_0}.}
\label{nuc12:eq:4.8}
\end{equation}



Ici, \(\mathcal V_b=V_c a_b^3\) est le volume physique au rebond. \(G\) et \(V_c\) s’annulent dans la première égalité ; \(E_0\) demeure comme énergie matérielle. Dans le secteur conservé \(\Lambda^3\mathbb C^4\), la proposition~\ref{vii:prop:operador-espin-car} démontre \(\widehat{\mathscr Q}_{\rm sp}=4I\), de sorte que \(q_\varrho=4\) pour tout état normalisé supporté dans ce secteur, et
\(\tau E_0/\hbar=|\gamma|/\sqrt{1+\gamma^2}\).
On obtient ainsi une échelle dynamique pour l’état spinoriel et la réalisation cosmologique spécifiés.


\subsection{Prolongement du transport de mémoire : réalisation quantique}

La réalisation quantique utilise le couplage réversible de la section~\ref{nuclear:10}. Nous considérons d’abord un plan canonique réel de matrice alternée \(\Omega=\left(\begin{smallmatrix}0&1\\-1&0\end{smallmatrix}\right)\). La partition nonadique détermine



\begin{equation}
Q=\frac13\begin{pmatrix}\sqrt8 I&-I\\I&\sqrt8 I\end{pmatrix},\qquad
U_H=Q^{\mathsf T}\operatorname{diag}(I,H)Q
=\frac19\begin{pmatrix}8I+H&\sqrt8(H-I)\\
\sqrt8(H-I)&I+8H\end{pmatrix}.
\label{nuc12:eq:5.1}
\end{equation}



L’opérateur \(H\in\operatorname{Sp}(2,\mathbb R)\) transporte ce plan. La matrice \(Q\) conserve le produit euclidien et la forme \(\Omega\oplus\Omega\) ; c’est pourquoi \(U_H\) est symplectique. La première composante définit le système observé et la seconde le registre de mémoire ; la transformation complète agit sur les deux.

La boucle ordonnée construite dans la section~\ref{nuclear:11} est



\begin{equation}
C=\begin{pmatrix}0&-1\\1&-1\end{pmatrix},\quad
P=\begin{pmatrix}1&1\\0&1\end{pmatrix},\quad
H_n=C^{-1}P^{-n}CP^n
=\begin{pmatrix}1+n&n^2\\n&n^2-n+1\end{pmatrix}.
\label{nuc12:eq:5.2}
\end{equation}



Pour chaque \(n\in\mathbb Z\), on a \(\det H_n=1\). En particulier, pour \(n=1\),


\begin{equation}
H_1=\begin{pmatrix}2&1\\1&1\end{pmatrix}
=F_1^2,\quad F_1=\begin{pmatrix}1&1\\1&0\end{pmatrix}
=P F_{\rm av}P^{-1}.
\label{nuc12:eq:5.3}
\end{equation}



Son spectre est \(\{\varphi^2,\varphi^{-2}\}\), par la reconnaissance ultérieure du mode d’or déjà engendré. L’entier \(n\) spécifie le nombre orienté d’itérations dans le lacet choisi. La collection des transports est déterminée algébriquement ; sa réalisation comme protocole physique conserve en outre les données de préparation et d’évolution correspondantes.

\subsubsection{État initial, transport et lecture}

La proposition~\ref{prop:ii-covarianza-gaussiana} construit la réalisation gaussienne de l’ellipse d’action. Dans la section positive \(\hbar_a>0\), avec \(A>0\) et \(|C^*/A|<1\), sa covariance est




\begin{equation}
V_\eta=\frac{\hbar_a}{2}M_\eta M_\eta^{\mathsf T},\qquad
M_\eta=\operatorname{diag}(e^{\eta/2},e^{-\eta/2}),\qquad
\eta=\operatorname{artanh}(C^*/A).
\label{nuc12:eq:5.4}
\end{equation}



Dans la représentation de Schrödinger sur \(L^2(\mathbb R,dQ)\), avec \(\widehat Q\) la multiplication par \(Q\) et \(\widehat P=-i\hbar_a\,d/dQ\), la relation \([\widehat Q,\widehat P]=i\hbar_a I\) vaut dans l’espace de Schwartz. L’état pur normalisé est



\[
\psi_\eta(Q)=(\pi\hbar_a e^\eta)^{-1/4}
\exp[-Q^2/(2\hbar_a e^\eta)].
\]



La préparation est le produit \(\psi_\eta\otimes\psi_\eta\). Le transport de \eqref{nuc12:eq:5.1} s’exprime dans la même coordinatisation de covariance par



\begin{equation}
\widetilde U=(M_\eta\oplus M_\eta)U_H
(M_\eta^{-1}\oplus M_\eta^{-1}).
\label{nuc12:eq:5.5}
\end{equation}



Le lemme~\ref{app:lem:implementacion-simplectica} fournit une implémentation unitaire de cette matrice symplectique à deux modes. Les implémentations qui diffèrent d’une phase globale produisent la même conjugaison de l’état. La covariance se transforme par \(V\mapsto\widetilde U V\widetilde U^{\mathsf T}\) ; l’état conjoint reste pur et sa trace partielle sur le second mode détermine la lecture visible, conformément à la construction de la section~\ref{app:gaussianas-espectro}.

\begin{theorem}[Mélange réduit par déformation relative]

La covariance de la lecture visible est



\begin{equation}
\boxed{V_{\rm vis}=
\frac{\hbar_a}{2}M_\eta
\frac{8I+HH^{\mathsf T}}9
M_\eta^{\mathsf T}.}
\label{nuc12:eq:5.6}
\end{equation}



Sa valeur symplectique normalisée est



\begin{equation}
\boxed{\nu(H):=\frac{2\sqrt{\det V_{\rm vis}}}{\hbar_a}
=\sqrt{1+\frac8{81}\left[\operatorname{tr}(HH^{\mathsf T})-2\right]}.}
\label{nuc12:eq:5.7}
\end{equation}
\end{theorem}



\begin{proof}
Soient \(T=(8I+H)/9\) et \(D=\sqrt8(H-I)/9\). Les deux modes initiaux sont indépendants et ont la même covariance. Dans la coordinatisation équilibrée de l’état initial, la covariance visible normalisée est \(TT^{\mathsf T}+DD^{\mathsf T}\). Le développement donne



\[
\begin{aligned}
TT^{\mathsf T}+DD^{\mathsf T}
&=\frac{64I+8(H+H^{\mathsf T})+HH^{\mathsf T}
+8[HH^{\mathsf T}-H-H^{\mathsf T}+I]}{81}\\
&=\frac{8I+HH^{\mathsf T}}9.
\end{aligned}
\]



Cela démontre \eqref{nuc12:eq:5.6}. Puisque \(\det M_\eta=1\) et \(\det(HH^{\mathsf T})=1\),
\(\det(8I+HH^{\mathsf T})=65+8\operatorname{tr}(HH^{\mathsf T})\).
La division par \(81\) démontre \eqref{nuc12:eq:5.7}. Les deux valeurs propres positives de \(HH^{\mathsf T}\) sont réciproques ; leur somme est au moins deux. En conséquence, \(\nu\geq1\), avec égalité exactement lorsque \(HH^{\mathsf T}=I\).
\end{proof}

L’échelle \(\hbar_a\) et la déformation \(\eta\) sont conservées jusqu’au calcul du déterminant. Leur simplification est covariante car l’état et l’opérateur sont transportés conjointement. L’expression \(HH^{\mathsf T}\) utilise le produit euclidien de la coordinatisation équilibrée de l’état initial ; un changement de coordinatisation transporte également ce produit métrique.

Un transport \(H\) orthogonal et symplectique peut avoir \(D\neq0\) tout en conservant le produit gaussien isotrope, avec \(\nu=1\). Le complément opératoire et l’intrication de l’état ont donc des critères distincts. Dans cette préparation, la déformation relative non isométrique caractérise exactement le cas \(\nu>1\).

\subsubsection{Évaluation exacte du lacet d’or}

Pour \eqref{nuc12:eq:5.2},



\begin{equation}
\operatorname{tr}(H_nH_n^{\mathsf T})-2
=n^2(2n^2-2n+5),\qquad
\nu_n^2=1+\frac8{81}n^2(2n^2-2n+5).
\label{nuc12:eq:5.8}
\end{equation}



L’identité résulte de la somme des carrés des quatre entrées. Pour \(n=1\), la somme est sept et



\begin{equation}
\boxed{\nu_1=\frac{11}{9},\quad
\operatorname{Tr}\rho_{\rm vis}^2=\frac9{11},\quad
\bar n_{\rm W}=\frac19,\quad
\frac{S_{\rm vis}}{k_B}=\frac{10}{9}\log10-\log9.}
\label{nuc12:eq:5.9}
\end{equation}



Le théorème~\ref{app:thm:espectro-gaussiano} démontre que l’état gaussien centré d’un mode de valeur symplectique \(\nu\) est unitairement équivalent à la densité géométrique de \eqref{app:eq:estado-geometrico}, dont l’occupation moyenne est \(\bar n_{\rm W}=(\nu-1)/2\). Pour \(\nu=11/9\), ses valeurs propres sont



\begin{equation}
\lambda_j=\frac9{10}\,10^{-j},\qquad j=0,1,2,\ldots.
\label{nuc12:eq:5.10}
\end{equation}



La série a pour somme un. Son carré a pour somme
\((9/10)^2/(1-10^{-2})=9/11\). Enfin,



\[
-\sum_j\lambda_j\log\lambda_j
=-\log(9/10)+\log10\sum_jj\lambda_j
=\frac{10}{9}\log10-\log9.
\]



La réduction gaussienne et son spectre sont démontrés dans la section~\ref{app:gaussianas-espectro} ; les deux séries précédentes évaluent exactement la pureté et l’entropie de cette réalisation. Comme l’état conjoint est pur, \(S_{\rm vis}/k_B\) est également son entropie adimensionnelle d’intrication. La transformation globale est unitaire, tandis que la restriction au premier mode a une entropie positive.

Le résultat correspond à la préparation et au transport spécifiés. Son identification à une mesure du vide requiert la réalisation expérimentale de ce protocole. L’entropie d’une réduction et la chaleur transférée restent des grandeurs distinctes.


\subsubsection{Composition avec le transducteur thermique}

Dans la représentation normale de Williamson, on choisit une fréquence \(\omega>0\) et on définit le lecteur énergétique
\(\widehat H_{\rm W}=\hbar_a\omega(\widehat N+\tfrac12)\). La fréquence fait partie de cette réalisation énergétique ; la covariance détermine les populations. Avec ce choix, \eqref{nuc12:eq:5.10} est un état de Gibbs dont le rapport des populations satisfait



\[
e^{-\hbar_a\omega/(k_BT_{\rm W})}=\frac1{10},\qquad
T_{\rm W}=\frac{\hbar_a\omega}{k_B\log10}.
\]



Si l’on prend la fréquence de l’horloge de la même section, \(\omega=1/t_P=2\pi/T_{108}\), l’identité du transducteur thermique \eqref{v:ant:eq:ii-kb-aut-par-termico},
\(\hbar_a/t_P=k_B\Theta_{{\rm clk},a}\log3\), donne




\begin{equation}
\boxed{\frac{T_{\rm W}}{\Theta_{{\rm clk},a}}=\frac{\log3}{\log10}.}
\label{nuc12:eq:5.11}
\end{equation}



La section d’action et le transducteur thermique s’annulent dans le quotient. L’indice \(a\) de \(\Theta_{{\rm clk},a}\) conserve la compatibilité avec la section \(\hbar_a\), comme dans \eqref{v:ant:eq:ii-kb-aut-orientaciones}. La température \(T_{\rm W}\) est l’équivalente de Gibbs pour le Hamiltonien déclaré ; l’évolution unitaire isolée conserve son caractère réversible. Cette identification spectrale détermine un quotient thermique et ne postule pas de dynamique de thermalisation.

\subsection{Extension à une dimension canonique finie arbitraire
}

Soit \(H\in\operatorname{Sp}(2d,\mathbb R)\), avec \(d\geq1\) fini. On prépare deux états gaussiens purs, centrés et identiques, de \(d\) modes chacun, de covariance \(V_0=(\hbar_a/2)MM^{\mathsf T}\), où \(M\in\operatorname{Sp}(2d,\mathbb R)\). La section~\ref{app:gaussianas-espectro} spécifie cette préparation et le transport conjoint du système de coordonnées par \(M\oplus M\). Dans le système de coordonnées équilibré, le même couplage \(U_H\) produit \(W=(8I+HH^{\mathsf T})/9\). Le produit \(A=HH^{\mathsf T}\) est positif, symétrique et symplectique.

Il existe une base orthonormale symplectique dans laquelle



\[
A=\operatorname{diag}(e^{2r_1},e^{-2r_1},\ldots,
e^{2r_d},e^{-2r_d}),\qquad r_j\ge0.
\]



L’identité \(A\Omega A=\Omega\), avec \(\Omega=\bigoplus_{j=1}^d\left(
\begin{smallmatrix}0&1\\-1&0\end{smallmatrix}\right)\), implique que \(\Omega\) envoie l’espace propre de \(\lambda\) sur celui de \(\lambda^{-1}\). On choisit des bases orthonormales appariées dans ces espaces. L’espace propre de \(\lambda=1\) est \(\Omega\)-invariant et admet des paires canoniques orthonormales. Leur réunion construit la base requise.


La matrice \(W\) est diagonale dans cette base. Les produits de chaque paire déterminent ses valeurs symplectiques :




\begin{equation}
\boxed{\nu_j^2=
\frac{(8+e^{2r_j})(8+e^{-2r_j})}{81}
=1+\frac{32}{81}\sinh^2r_j.}
\label{nuc12:eq:6.1}
\end{equation}



Le théorème~\ref{app:thm:espectro-gaussiano} compose la réduction par trace partielle avec ces valeurs symplectiques :



\begin{equation}
\boxed{\operatorname{Tr}\rho_{\rm vis}^2=\prod_{j=1}^d\nu_j^{-1},
\qquad
\frac{S_{\rm vis}}{k_B}=\sum_{j=1}^d
\left[\frac{\nu_j+1}{2}\log\frac{\nu_j+1}{2}
-\frac{\nu_j-1}{2}\log\frac{\nu_j-1}{2}\right].}
\label{nuc12:eq:6.2}
\end{equation}



On adopte \(0\log0=0\) par continuité. La preuve du théorème~\ref{app:thm:espectro-gaussiano} transforme unitairement l’état en un produit de \(d\) densités géométriques : la pureté du produit est le produit des puretés et son entropie est la somme des entropies. Le résultat vaut pour chaque \(d\) fini et chaque transport \(H\) du domaine, en maintenant toutes les occupations de l’espace de Fock.

Dans la boucle \eqref{nuc12:eq:5.3}, \(r=2\log\varphi\) et \(\sinh^2r=5/4\), ce qui reproduit \eqref{nuc12:eq:5.9}. Ainsi, le spectre du transport et la partition nonadique \(8+1\) déterminent le spectre de la lecture réduite ; les sections d’action fixent sa réalisation dimensionnelle.

La dimension d’une réalisation exceptionnelle et son transport proviennent du constructeur correspondant. L’équation \eqref{nuc12:eq:6.1} détermine la loi spectrale après la réalisation de ses modes canoniques et de la préparation gaussienne spécifiée.


\subsubsection{Identité opératorielle entre mémoire et structure complexe
}

Dans le système de coordonnées canonique équilibré, soit \(J=\Omega\), avec \(J^2=-I\) et \(J^{\mathsf T}=-J\). Cette matrice réalise la structure complexe canonique de l’espace des phases, et \(D=\sqrt8(H-I)/9\) est le bloc de mémoire. Pour la même préparation et le même transport de système de coordonnées de la section~\ref{app:gaussianas-espectro}, on a



\begin{equation}
\boxed{-(JW)^2=I+[D,J][D,J]^{\mathsf T},\qquad
W=\frac{2}{\hbar_a}M^{-1}V_{\rm vis}M^{-\mathsf T}.}
\label{nuc12:eq:6.3}
\end{equation}



\begin{proof}
Les identités symplectiques \(HJH^{\mathsf T}=J\) et \(H^{\mathsf T}JH=J\), avec \(A=HH^{\mathsf T}\), donnent \(JAJ^{\mathsf T}=A^{-1}\). Si \(C=HJ-JH\), la multiplication des quatre termes produit




\[
CC^{\mathsf T}
=A+JAJ^{\mathsf T}
-HJH^{\mathsf T}J^{\mathsf T}
-JHJ^{\mathsf T}H^{\mathsf T}
=A+A^{-1}-2I.
\]



D'autre part,



\[
-(JW)^2
=\frac{(8I+A^{-1})(8I+A)}{81}
=I+\frac8{81}(A+A^{-1}-2I).
\]



Comme \([D,J]=\sqrt8\,C/9\), on obtient \eqref{nuc12:eq:6.3}. L’identité vaut dans chaque dimension canonique finie et pour chaque \(H\) du domaine, y compris les transports qui mélangent les modes.
\end{proof}

Les valeurs propres de cet opérateur sont les \(\nu_j^2\), chacune de multiplicité deux. La formule~\eqref{app:eq:pureza-general} exprime alors la pureté directement en termes du bloc de mémoire :




\begin{equation}
\boxed{\operatorname{Tr}\rho_{\rm vis}^2
=\det\!\left(I+[D,J][D,J]^{\mathsf T}\right)^{-1/4}.}
\label{nuc12:eq:6.4}
\end{equation}



Dans un mode,



\begin{equation}
\nu^2=1+\frac12\|[D,J]\|_F^2.
\label{nuc12:eq:6.5}
\end{equation}



Dans ce protocole, \(D\) commute avec \(J\) exactement lorsque la lecture réduite conserve la pureté. Pour la boucle \(H_1\), \(\|[D,J]\|_F^2=80/81\), et \eqref{nuc12:eq:6.5} donne \(\nu^2=121/81\), d’où la pureté \(9/11\).

Le commutateur utilise la structure complexe compatible avec la covariance initiale et son produit métrique. Un changement canonique transporte conjointement les trois données ; la transposée de \eqref{nuc12:eq:6.3} se réfère au système de coordonnées équilibré. La construction préalable de l’élément imaginaire dans le noyau formel précède cette représentation, qui spécifie son action canonique.

L’évolution conjointe conserve la pureté par sa mise en œuvre unitaire ; le commutateur détermine la perte de pureté de la lecture partielle. L’information de l’état conjoint demeure dans les deux facteurs et dans leurs corrélations. La formule détermine la pureté de la réduction, qui est une fonctionnelle non linéaire de son opérateur densité.


La décomposition réelle par rapport à \(J\) distingue ses composantes complexe linéaire et complexe antilinéaire :



\[
D_{\rm lin}=\frac{D-JDJ}{2},\qquad
D_{\rm anti}=\frac{D+JDJ}{2}.
\]



La première composante commute avec \(J\) et la seconde anticommute avec \(J\). Les identités
\([D,J]=2D_{\rm anti}J\) et \(JJ^{\mathsf T}=I\) impliquent



\begin{equation}
\boxed{-(JW)^2=I+4D_{\rm anti}D_{\rm anti}^{\mathsf T},\qquad
\operatorname{Tr}\rho_{\rm vis}^2
=\det(I+4D_{\rm anti}D_{\rm anti}^{\mathsf T})^{-1/4}.}
\label{nuc12:eq:6.6}
\end{equation}



La composante complexe antilinéaire du bloc de mémoire, définie par rapport à la structure complexe de l’état initial, détermine exactement la pureté et le spectre de mélange de ce protocole. La composante complexe linéaire peut transporter la mémoire en conservant la pureté. La distinction se rapporte à la linéarité par rapport à \(J\) ; la complexité combinatoire et la classification arithmétique des constantes sont des propriétés d’autres objets.
